\chapter*{Abstract}

Promoter region identification is an important step in functional genome
annotation because promoters regulate the initiation of gene transcription. This
thesis investigates machine learning approaches for promoter prediction in DNA
sequences, with emphasis on capsule-based architectures and representations
derived from physicochemical properties. The research builds on the qualification
stage, including the CapsProm model and the physicochemical-feature study, and
reorganizes these contributions into a reproducible experimental protocol. The
work considers source-defined datasets derived from CNNProm, GraphPro, Bhandari et
al., and CDBProm, normalized to 81 bp windows around the transcription start
site. Experiments will be executed with explicit provenance control, formal experiment specification, and thesis-ready results generated from the reproducible protocol.

\noindent \textbf{Keywords:} bioinformatics; promoter prediction; capsule neural
networks; machine learning; DNA sequences.
